\section{Original input and physical-output resources}
\label{sec:peps_physical_first_resources}

The following bounds concern the actual source replacement and its prescribed
star and sample alphabets. They do not assert an identification of the sampled
physical operator with a contraction of local tensors.

\begin{theorem}[Finite coordinates for the original input]
    \label{thm:peps_finite_whole_party_input}
    \lean{TNLean.PEPS.PairEffect.ProductInput.finiteDimensional_wholePartyMem}
    \leanok
    Let $P$ be a finite set of parties and let $H_p$ be their original
    finite-dimensional private memories. The ordered memory
    $\bigotimes_{p\in P}H_p$ is finite-dimensional and hence admits a finite
    orthonormal basis. There is no numerical bound on the dimensions $\dim H_p$.
\end{theorem}
\begin{proof}
    \leanok
    Each register in the original party list has space $H_p$ for one $p\in P$.
    Finite tensor products of finite-dimensional spaces are finite-dimensional.
    This includes the empty tensor product and zero-dimensional factors.
\end{proof}


\begin{theorem}[Source positions after the physical output exchange]
    \label{thm:peps_physical_first_source_count}
    \lean{TNLean.PEPS.PairEffect.OriginalCircuit.card_sourceLocations_physicalFirstReplacement_le}
    \leanok
    Let $W$ be an original contraction circuit with $N$ nonprivate gate
    occurrences, each involving at most $a$ parties. Its physical-first
    source replacement has at most $N\binom a2$ source positions. This count
    does not depend on the effect-elimination accuracy or on private dimensions.
\end{theorem}
\begin{proof}
    \uses{thm:peps_physical_first_occurrences}
    \leanok
    Each replacement gate has one source slot per unordered pair of its
    original participants. The final physical output exchange introduces
    no new gate or source occurrence.
\end{proof}

\begin{theorem}[Polynomial sample count from original resources]
    \label{thm:peps_physical_first_sample_polynomial}
    \lean{TNLean.PEPS.PairEffect.OriginalCircuit.sourceSamplingCount_physicalFirstReplacement_lt_power}
    \leanok
    Let $W,N,a$ be as in the preceding source-count theorem.
    Let $L\ge1$, let $n,s,d,p$ be nonnegative integers, and suppose
    \begin{align}
        N\le C_NL^n,\qquad 0\le B\le C_BL^s,\qquad 0\le D\le C_DL^d,
        \qquad C_B\ge0.
    \end{align}
    Here $B$ bounds the original gate coefficient sums, $D$ is a nonnegative
    parameter, and $\ell$ is a nonnegative integer. In applications they are
    chosen as a physical-dimension bound and a whole-lifetime participation bound.
    Put $C_{\rm src}=\max\{1,C_N\binom a2\}$.
    For the actual sample count of the physical-first replacement, chosen at
    accuracy $L^{-p}$ and cost $B^{4\ell}D^4$, one has
    \begin{align}
        k<\bigl[64C_{\rm src}^2(C_B^{4\ell}C_D^4)^2+2\bigr]
        L^{2(n+4\ell s+4d+p)}.
    \end{align}
    Every coefficient and exponent in this bound depends only on the fixed
    original resource bounds and the prescribed accuracy.
\end{theorem}
\begin{proof}
    \leanok
    \uses{thm:peps_physical_first_source_count,thm:peps_source_sampling_polynomial}
    The source count is at most $C_{\rm src}L^n$. Apply the explicit
    sample-count estimate with this bound and the displayed bounds on $B,D$.
    The maximum with one covers the case of no sources.
\end{proof}

\begin{theorem}[Polynomial branch counts after the physical exchange]
    \label{thm:peps_physical_first_branch_polynomial}
    \lean{TNLean.PEPS.PairEffect.OriginalCircuit.card_branchLabels_physicalFirstReplacement_le_power}
    \leanok
    Suppose $L\ge1$, $C_N\ge1$, $C_K,C_S,S\ge0$, and
    \begin{align}
        K\le C_KL^u,\qquad N\le C_NL^n,\qquad S\le C_SL^s.
    \end{align}
    Let every original gate contain at most $K$ monomials with at most $r$
    pair effects each and coefficient sum at most $S$. Use gate budget
    $\delta=L^{-p}/(8\max\{1,N\})$, where $u,n,s,p$ are nonnegative integers.
    Every gate of the actual physical-first replacement then has at most
    \begin{align}
        C_K\bigl[64C_N^2(rC_S)^2+2\bigr]^r
        L^{u+2r(n+s+p)}
    \end{align}
    branches.
\end{theorem}
\begin{proof}
    \leanok
    \uses{thm:peps_effect_replacement_polynomial}
    The replacement has at most $K m^r$ branches, where $m$ is the chosen
    stack length. Substitute its polynomial bound. The physical output
    exchange changes neither branch labels nor coefficient sums.
\end{proof}

\begin{definition}[Original gate occurrences]
    \label{def:peps_physical_first_original_locality}
    \lean{TNLean.PEPS.PairEffect.OriginalCircuit.physicalFirstReplacementLocationsEquiv}
    \leanok
    \uses{thm:peps_replacement_participants,thm:peps_physical_first_occurrences}
    Let $E$ be the bijection from original nonprivate gate occurrences to
    physical-first replacement gate occurrences obtained by composing the
    correspondences for effect elimination and the physical output exchange.
\end{definition}

\begin{theorem}[Original participants and link locality]
    \label{thm:peps_physical_first_original_locality}
    \lean{TNLean.PEPS.PairEffect.OriginalCircuit.participants_physicalFirstReplacementLocationsEquiv}
    \lean{TNLean.PEPS.PairEffect.OriginalCircuit.exists_original_gate_of_physicalFirst_link}
    \uses{def:peps_physical_first_original_locality}
    \leanok
    The bijection $E$ preserves participating sets. Choose any collection of
    replacement gate occurrences and a root assignment on all replacement
    occurrences whose value participates at every retained occurrence.
    Every star or sample link in the resulting construction joins two parties
    that participate in a common gate of the original circuit.
\end{theorem}
\begin{proof}
    \leanok
    \uses{thm:peps_replacement_participants,thm:peps_physical_first_occurrences,
        thm:compression_original_gate_locality}
    The two gate correspondences preserve participating sets, so their
    composition does as well. A link's gate label places both endpoints in
    its replacement gate's participating set; the correspondence identifies
    this set with the original one.
\end{proof}

\begin{theorem}[Original lifetime bound controls the actual link degree]
    \label{thm:peps_physical_first_original_degree}
    \lean{TNLean.PEPS.PairEffect.OriginalCircuit.card_incidentPhysicalFirstLinks_le}
    \leanok
    Suppose every original gate involves at most $b$ parties and each party
    participates in at most $b$ original nonprivate occurrences over the
    whole construction. For the retained gate collection and participating
    roots above, each party is incident to at most $2b(b-1)$ actual links,
    counting parallel links separately.
\end{theorem}
\begin{proof}
    \leanok
    \uses{thm:peps_physical_first_original_locality}
    The gate correspondence preserves arity and lifetime participation.
    Discarding some gate occurrences can only reduce the latter count.
    Each participating gate contributes at most $b-1$ star links and $b-1$
    sample links at a party. Sum over its at most $b$ original occurrences.
\end{proof}

\begin{theorem}[Polynomial virtual dimensions of the actual source alphabets]
    \label{thm:peps_source_alphabet_polynomial}
    \lean{TNLean.PEPS.PairEffect.SourceCircuit.card_sourceNetworkAlphabet_le_of_power_bounds}
    \leanok
    Fix a set of gate occurrences of an actual source circuit, and the
    root assignment defining its star links. A star link at occurrence
    $g$ carries the ordered ket--bra pair of its actual branch labels;
    each sample link carries one of $k$ sample indices.
    Suppose $L\geq1$, $C_M,C_k\geq0$, and $u,v$ are nonnegative integers.
    If each of these gate occurrences has at most $C_ML^u$ branches and
    $k\leq C_kL^v$, then every constructed link alphabet satisfies
    \begin{align}
        \dim(\text{link alphabet})
        \leq (C_M^2+C_k)L^{\max(2u,v)}.
        \label{eq:peps_source_alphabet_polynomial}
    \end{align}
    No dimension of a private memory or source Schmidt space occurs in
    this bound. The statement concerns the actual alphabets; the
    identification of the sampled physical operator with a contraction
    of local tensors is a separate assertion.
\end{theorem}
\begin{proof}
    \leanok
    At a star link the cardinality is the square of the actual branch
    count, hence at most $C_M^2L^{2u}$. At a sample link it is exactly
    $k\leq C_kL^v$. Since $L\geq1$, increase either exponent to
    $\max(2u,v)$ and bound its nonnegative coefficient by $C_M^2+C_k$.
    % Source: 04-compression.tex, lines 565--588.
    These are the star and sample alphabets in
    \cite[Proof of Theorem~5.2]{OpenAI2026PolynomialPEPS}.
\end{proof}
